\documentclass{article}

\usepackage{iclr2025_conference,times}
\usepackage{hyperref}
\usepackage{url}
\usepackage{amsmath,amssymb}
\usepackage{booktabs}
\usepackage{graphicx}
\usepackage{subcaption}
\usepackage{microtype}
\usepackage{enumitem}

\title{Continual Learning on Split-CIFAR-100:\\
Overcoming Catastrophic Forgetting Without Replay}

\author{%
  \parbox[t]{0.4\linewidth}{\centering
    Divya Sree Durga Devalla \\
    \texttt{ddeva@uic.edu}
  }%
  \hfill
  \parbox[t]{0.4\linewidth}{\centering
    Yuvaneswaren Ramakrishnan Sureshbabu \\
    \texttt{yrama@uic.edu}
  }\\[8pt]
  \centering
  Department of Computer Science \\
  University of Illinois Chicago
}

\iclrfinalcopy

\begin{document}

\maketitle

% ================================================================
\begin{abstract}
Catastrophic forgetting is one of the most persistent problems in sequential
neural network training. When a model is updated on a new task, gradient updates
tend to overwrite the weights that encoded earlier tasks, causing prior knowledge
to simply disappear. This project explores how far a model can get on
Split-CIFAR-100, a standard five-task continual learning benchmark, without
storing or replaying any past training data. We implement and compare eight
methods, starting from naive fine-tuning and working toward architectural
solutions that sidestep forgetting rather than fight it with penalties. Our best
configuration, a frozen Vision Transformer (ViT-B/16) backbone combined with
per-task adapter networks and a non-parametric Nearest Class Mean (NCM)
classifier, achieves 74.81\% average accuracy with only $-9.64\%$ backward
transfer. That is a roughly six-fold improvement in accuracy and a seven-fold
reduction in forgetting compared to the naive baseline (12.86\% AA,
$-60.95\%$ BWT). The central finding is that eliminating trainable classifier
weights and expanding the feature space incrementally per task is far more
effective than any regularization scheme applied to a shared, mutable network.
\end{abstract}

% ================================================================
\section{Introduction}
\label{sec:intro}

\subsection{The Problem}

Humans accumulate knowledge over a lifetime without earlier skills disappearing
when new ones are acquired. Neural networks trained with standard gradient
descent do not share this property. When a network is updated on a new task,
the parameter updates that improve performance on that task tend to overwrite
whatever the network encoded before. This phenomenon, known as
\textbf{catastrophic forgetting} \citep{kirkpatrick2017overcoming}, causes a
model to excel at the most recent task while performing near chance on everything
that came before it.

This limitation has real consequences outside the lab. A radiology model that
forgets earlier disease categories each time it is updated on a new disease
cannot be trusted clinically. A robot controller that loses navigation skills
whenever it learns a new manipulation behavior is impractical for deployment.
The field of \textbf{continual learning} (also called lifelong or incremental
learning) is dedicated to solving this problem: enabling a model to absorb new
tasks sequentially without losing its grip on earlier ones \citep{li2017learning}.

\subsection{Previous Work}

Proposed solutions broadly fall into three families. \textbf{Replay-based}
methods store a buffer of past examples and mix them into each new training
round \citep{zhu2021prototype}. While effective, they raise storage and privacy
concerns and are often impractical when past data simply cannot be kept.
\textbf{Regularization-based} methods add penalty terms to the training loss
that discourage large changes to parameters important for earlier tasks. Elastic
Weight Consolidation (EWC) \citep{kirkpatrick2017overcoming} uses the Fisher
Information Matrix to identify and protect critical weights, while Learning
Without Forgetting (LwF) \citep{li2017learning} uses knowledge distillation to
constrain output-space drift. \textbf{Architecture-based} methods take a
different angle: rather than constraining how a shared network changes, they
allocate dedicated capacity to each task so that new learning never touches old
representations. Recent work has shown that combining strong pretrained
backbones with task-specific modules can dramatically reduce forgetting in
class-incremental settings \citep{zhou2024revisiting, zhang2023slca,
zhu2021prototype}.

A recurring theme in recent literature is that the quality of the pretrained
backbone matters as much as the anti-forgetting strategy. Models pretrained on
large datasets produce feature spaces rich enough that simple non-parametric
classifiers can achieve strong performance, which sidesteps the classifier
weight overwriting problem entirely \citep{zhou2024revisiting}.

\subsection{Our Methods}

This project takes a strictly replay-free stance: no past data is stored at any
point. Within that constraint, we implemented and compared eight methods in order
of increasing sophistication. The first four use a pretrained ResNet-18 backbone
with a growing linear classification head and progressively stronger
regularization: naive fine-tuning, Learning Without Forgetting, a hybrid of EWC
and LwF, and standalone EWC. The fifth group repeats these regularization methods
on a stronger frozen ViT-B/16 backbone. The final three methods abandon the
trainable classification head entirely, instead using non-parametric Nearest
Class Mean (NCM) inference over class prototypes stored in a fixed or
incrementally expanding feature space. The best of these pairs a frozen
ViT-B/16 backbone with per-task MLP adapters and NCM inference.

\subsection{Summary of Results}

The experiments tell a clear story. Regularization-based methods on ResNet-18
plateau below 22\% average accuracy, regardless of which penalty is used.
Switching to a stronger ViT-B/16 backbone nearly doubles this, but forgetting
in the shared classifier remains severe. The decisive change comes from replacing
the trainable head with an NCM classifier: this single architectural shift cuts
backward transfer from roughly $-40\%$ to roughly $-10\%$ across all backbone
types. Combining a frozen ViT-B/16 backbone with per-task adapters and NCM
reaches 74.81\% average accuracy and $-9.64\%$ backward transfer, without
storing a single past training example and without needing task identity at
test time.

% ================================================================
\section{Problem Description}
\label{sec:problem}

\subsection{Task Setup and Evaluation}

We work in the \textbf{class-incremental} continual learning setting, which is
the more challenging of the standard continual learning protocols.

\textbf{Split-CIFAR-100.} The primary benchmark partitions the 100 classes
of CIFAR-100 \citep{zenke2017continual} into five disjoint groups of 20 classes
each. The model is trained sequentially on Tasks 1 through 5 (classes 0--19,
20--39, 40--59, 60--79, 80--99). At training time, only the current task's data
is available and no replay buffer is maintained. At test time, the model must
classify a query among all classes seen so far without any task identity
information being provided. By the final task, the model is distinguishing among
all 100 classes simultaneously.

\textbf{CORe50.} To test whether our best approach generalises beyond
Split-CIFAR-100, we additionally evaluate on CORe50 \citep{lomonaco2017core50},
a continual learning benchmark built around 50 object categories captured across
11 different sessions (varying backgrounds, lighting, and viewpoints). 

The two metrics tracked throughout are:
\begin{itemize}[leftmargin=*]
  \item \textbf{Average Accuracy (AA):} the mean of per-task test accuracies
    evaluated after all five tasks have been trained. This captures both
    plasticity (learning new tasks well) and stability (retaining old ones).
  \item \textbf{Backward Transfer (BWT):} the average drop in accuracy on
    previously seen tasks caused by training on subsequent ones. A BWT of 0
    means no forgetting; more negative values indicate more severe forgetting.
\end{itemize}

Formally, let $R_{i,j}$ denote test accuracy on Task $i$ right after training
on Task $j$, where $j \geq i$. Then:
\begin{align}
  \text{AA}  &= \frac{1}{T} \sum_{i=1}^{T} R_{i,T}, \\[4pt]
  \text{BWT} &= \frac{1}{T-1} \sum_{i=1}^{T-1}
               \bigl(R_{i,T} - R_{i,i}\bigr),
\end{align}
where $T = 5$ is the total number of tasks.

% -------------------------------------------------------
\subsection{Methods}
\label{sec:methods}

We implemented eight methods in total, each building on the lessons of the
previous one. All experiments use PyTorch with Adam optimizer ($lr = 10^{-3}$)
and 5 epochs per task unless otherwise noted.

The core problem with sequential training is that every gradient update on a
new task pushes the model's weights in a direction that improves new-task
performance, but nothing forces those updates to stay within the region of
weight space that preserved old-task performance. Old knowledge is simply
overwritten. The methods below attack this problem with increasing architectural
ambition, from adding soft penalties on top of the same shared network, to
abandoning the shared network structure altogether.

\subsubsection{Naive Fine-Tuning (Baseline)}

The simplest approach: train on each task in sequence with no modifications,
using a pretrained ResNet-18 backbone with a growing linear classification head
that adds 20 output neurons per task. Nothing prevents forgetting, making this
the lower bound all other methods must beat.
As expected, the model forgets nearly everything it has learned by the final task, with accuracy collapsing to near chance across all earlier tasks.

\subsubsection{Learning Without Forgetting (LwF)}

Before each new task, the current model is frozen as a ``teacher.'' As the
live model trains on new data, it is also penalised for drifting away from the
teacher's predictions on old classes. The distillation loss is:
\begin{equation}
  \mathcal{L}(\theta) \;=\; \mathcal{L}_{\mathrm{CE}}(\theta) \;+\;
  \lambda_{\mathrm{KD}} \cdot \mathcal{L}_{\mathrm{KD}}(\theta),
  \label{eq:lwf_total}
\end{equation}
where $\lambda_{\mathrm{KD}}$ grows with each task. The limitation is that
the teacher only sees new-task images, so its guidance weakens as the backbone
shifts over time.
Distillation provides a real but modest improvement over naive fine-tuning, and forgetting remains severe because the teacher's guidance degrades as the backbone shifts.

\subsubsection{Elastic Weight Consolidation (EWC)}

Instead of constraining predictions, EWC \citep{kirkpatrick2017overcoming}
constrains which weights are allowed to change. After each task, the Fisher
Information Matrix identifies which weights mattered most. Those weights are
then anchored by a quadratic penalty during the next task:
\begin{equation}
  \mathcal{L}(\theta) \;=\; \mathcal{L}_{t+1}(\theta) \;+\;
  \frac{\lambda}{2} \sum_{k} F^{(t)}_k
  \bigl(\theta_k - \theta_{t,k}^*\bigr)^2,
  \label{eq:ewc}
\end{equation}
where $\lambda = 10{,}000$. Important weights are held in place; unimportant
ones remain free to adapt.
EWC is the strongest regularization method tested, though it still plateaus well below what architectural changes later achieve.

\subsubsection{Hybrid (EWC + LwF)}

Both penalties are applied simultaneously: EWC protects important weights
while LwF keeps output predictions consistent with the teacher.
\begin{equation}
  \mathcal{L}(\theta) \;=\;
  \mathcal{L}_{\mathrm{CE}} \;+\;
  \lambda_{\mathrm{EWC}} \cdot \mathcal{L}_{\mathrm{EWC}} \;+\;
  \lambda_{\mathrm{KD}}^{\mathrm{adapt}} \cdot \mathcal{L}_{\mathrm{KD}}.
\end{equation}
In practice the two constraints can pull in opposite directions, limiting
how much each individually contributes.
Combining the two penalties improves backward transfer slightly over standalone LwF, but adds little over EWC alone since the two objectives partially conflict.

Across all four regularization methods, replacing linear head weights with L2-normalised class mean vectors after each task (prototype alignment correction) reduces backward transfer from roughly $-39\%$ to $-9\%$ with no loss in average accuracy, because a frozen backbone keeps the feature space stable between tasks. Adding backbone fine-tuning raises per-task training accuracy to $59$--$63\%$ but reintroduces catastrophic forgetting, since shifting the feature space for a new task invalidates the stored prototypes of all prior ones. The best configuration within the regularization family is therefore a frozen backbone with alignment correction, and even this reaches a hard ceiling that the architectural changes below later break through.

\subsubsection{ViT-B/16 Backbone with Regularization}

Rather than changing the training strategy, we ask: what if the model starts
with better features? We replace ResNet-18 with a frozen ViT-B/16
\citep{dosovitskiy2020image}, raising feature dimensionality from 512 to 768.
Only the classification head is trained. The shared mutable head problem
remains, but stronger features raise the accuracy ceiling noticeably.
Better features lift accuracy substantially across all three methods, but forgetting in the shared classifier remains just as persistent as with ResNet-18.

\subsubsection{Prototype-Based Classification (PASS)}

All methods above share one structural weakness: a single set of classification
weights that all tasks must overwrite. Instead, we store the average feature
vector of each class as a \textbf{prototype} and classify by nearest distance:
\begin{equation}
  \mu_c = \frac{1}{|S_c|} \sum_{x \in S_c} f_\theta(x), \qquad
  \hat{y} = \arg\min_{c}\;\|f_\theta(x) - \mu_c\|_2.
\end{equation}
With a frozen backbone, no prototype can interfere with any other; forgetting
in the classifier is structurally impossible. Each class is also augmented with
50 synthetic feature vectors sampled from its feature covariance to produce
more robust prototypes.
Switching to a prototype-based classifier more than doubles accuracy over the best regularization result, with far less forgetting and no penalty terms at all.

\subsubsection{ADAM: Per-Task Adapters with NCM (Our Best Method)}

Even with frozen features, packing 100 class prototypes into a fixed-dimensional
space causes crowding. ADAM solves this by giving each task its own small
two-layer MLP \textbf{adapter} that is trained on that task and then permanently
frozen. After $t$ tasks, the final feature is the original backbone output
concatenated with all adapter outputs, giving an expanding
$(d + t \times 128)$-dimensional vector where $d$ is 512 for ResNet-18 and
768 for ViT-B/16. Class prototypes are recomputed in this growing space after
every task, and a Tukey power transform with L2 normalisation is applied
before NCM inference. New tasks add capacity without touching anything
already learned.
This is our best result: high accuracy across all tasks on both benchmarks, with near-zero forgetting and no stored training data.
Figure~\ref{fig:adam_flowchart} illustrates the full pipeline.

\begin{figure}[h]
  \centering
  \includegraphics[width=0.55\linewidth]{adam_flowchart.png}
  \caption{ADAM architecture: frozen backbone features are concatenated with
    per-task adapter outputs, Tukey-transformed, and classified via NCM.}
  \label{fig:adam_flowchart}
\end{figure}

% ================================================================
\section{Results}
\label{sec:results}

Table~\ref{tab:results} summarizes average accuracy and backward transfer for
all ten configurations. Figure~\ref{fig:comparison} shows per-task forgetting
curves for three representative configurations spanning the full range of
behaviours observed. Figure~\ref{fig:ablation_bars} details the complete
ablation over all four regularization methods and three training regimes.

\begin{table}[h]
\centering
\caption{All methods on Split-CIFAR-100 in the class-incremental setting.
AA = Average Accuracy after all five tasks. BWT = Backward Transfer (closer
to 0 is better). All methods are fully replay-free.}
\label{tab:results}
\renewcommand{\arraystretch}{1.25}
\begin{tabular}{llccc}
\toprule
\textbf{Method} & \textbf{Backbone} & \textbf{Classifier}
  & \textbf{AA (\%)} & \textbf{BWT (\%)} \\
\midrule
Naive Fine-Tuning                  & ResNet-18 & Linear Head & 12.86 & $-60.95$ \\
LwF                                & ResNet-18 & Linear Head & 16.58 & $-56.09$ \\
Hybrid (EWC + LwF)                 & ResNet-18 & Linear Head & 17.56 & $-43.24$ \\
EWC                                & ResNet-18 & Linear Head & 21.73 & $-40.66$ \\
\midrule
ViT Naive                          & ViT-B/16  & Linear Head & 36.32 & $-69.58$ \\
ViT LwF                            & ViT-B/16  & Linear Head & 36.15 & $-69.62$ \\
ViT EWC                            & ViT-B/16  & Linear Head & 41.14 & $-62.88$ \\
\midrule
ResNet-18 + Adapters + NCM         & ResNet-18 & NCM         & 47.63 & $-15.54$ \\
Frozen ResNet-18 + Proto Aug + NCM & ResNet-18 & NCM         & 57.20 & $-10.79$ \\
\textbf{ViT + Adapters + NCM}
  & \textbf{ViT-B/16} & \textbf{NCM}
  & \textbf{74.81} & $\mathbf{-9.64}$ \\
\bottomrule
\end{tabular}
\end{table}

\begin{figure}[h]
  \centering
  \includegraphics[width=\linewidth]{final_fig3_forgetting_curves.png}
  \caption{Per-task accuracy over the training sequence for three representative
    configurations: Naive (complete collapse), EWC+Align (gradual decay, retained
    memory), and LwF+Align+FT (high plasticity but sharp forgetting).}
  \label{fig:comparison}
\end{figure}

\begin{figure}[h]
  \centering
  \includegraphics[scale=0.4]{final_fig1_grouped_bars.png}
  \caption{AA (top) and BWT (bottom) for all twelve regularization configurations
    (four methods $\times$ three regimes: Base, +Align, +Align+FT). Alignment
    correction alone cuts BWT from $-39\%$ to $-9\%$; backbone fine-tuning reverses this gain.}
  \label{fig:ablation_bars}
\end{figure}

\subsection{Regularization-Based Methods on ResNet-18}

The naive baseline makes the severity of catastrophic forgetting immediately
visible. Each task reaches around 60\% accuracy right after its own training,
then collapses toward 0\% as later tasks are trained
(Figure~\ref{fig:comparison}). The peak single-task accuracy of 63.9\% stands
in stark contrast to the final AA of just 12.86\%, and by the time Task 5 is
complete the model has essentially forgotten everything about Tasks 1 through 4,
with a BWT of $-60.95\%$.

LwF improves AA to 16.58\% and BWT to $-56.09\%$, a real but modest gain. Task
1 still drops to near 0\% by the end of training. The distillation objective
constrains the model's output behavior, but it operates only on new-task inputs,
making it an imperfect proxy for preserving old-task knowledge. When layer 4 is
also fine-tuned, backbone feature drift weakens the old teacher signals further
over time.

The Hybrid EWC and LwF combination reaches 17.56\% AA and $-43.24\%$ BWT.
Combining the two objectives helps BWT more than AA. In practice the two
regularizers work against each other to some degree: EWC anchors specific
parameter values while LwF simultaneously pushes output distributions to track
a frozen snapshot. These competing gradients reduce how much each penalty can
individually contribute.

Standalone EWC is the best of the four regularization methods on ResNet-18,
achieving 21.73\% AA and $-40.66\%$ BWT. The Fisher-weighted penalty is more
targeted than LwF's behavioral constraint because it works directly on parameter
values ranked by importance, rather than approximating importance through output
space. Even so, 21.73\% across five tasks reveals a hard ceiling. The cause is
structural: all tasks share the same weight space, and regularization can only
slow drift, not eliminate the capacity conflict.

\subsection{Stronger Features: ViT-B/16 with a Linear Head}

Replacing ResNet-18 with a frozen ViT-B/16 backbone nearly doubles accuracy
across equivalent methods. ViT EWC reaches 41.14\%, compared to ResNet-18 EWC
at 21.73\%. The richer 768-dimensional features give the linear classifier
significantly better material to work with from the start.

Interestingly, ViT Naive and ViT LwF perform almost identically (36.32\% vs.
36.15\%) and both show worse BWT than their ResNet-18 counterparts. With the
backbone frozen, only the growing classification head is trainable. LwF's
distillation over old head outputs is fighting against a head that expands by 20
neurons at every task in a class-incremental setting where all accumulated
classes must be distinguished simultaneously. Better features raise the accuracy
ceiling, but the forgetting mechanism in the shared classifier is completely
intact.

\subsection{Structural Fix: NCM Classifier with a Frozen Backbone}

The biggest single improvement in the entire experiment series happens when the
trainable head is replaced with NCM. The Frozen ResNet-18 with Prototype
Augmentation and NCM method achieves 57.20\% AA and $-10.79\%$ BWT using no
regularization at all. Compared to the best regularization result (ResNet-18
EWC: 21.73\% AA, $-40.66\%$ BWT), this is a 2.6x gain in accuracy and a 4x
reduction in forgetting, from a change in architecture alone.

Task 1's accuracy trajectory tells the story clearly. Under naive fine-tuning,
Task 1 goes from roughly 60\% right after training to nearly 0\% after five
tasks: near-total collapse. Under Frozen ResNet-18 with NCM, Task 1 goes from
78.7\% to 58.2\% across the same five tasks. This is real degradation but it is
gradual, caused by feature space crowding as 100 class prototypes are packed
into a 512-dimensional space, not by weight overwriting.

Prototype augmentation (sampling 50 synthetic feature vectors per class from
the class-level feature covariance) helps reduce crowding by giving each class
a more statistically robust centroid estimate, particularly for high-variance
classes.

\subsection{Per-Task Adapters on ResNet-18}

The ResNet-18 with Per-Task Adapters and NCM method adds a dedicated MLP adapter
per task. Somewhat counterintuitively, it scores below the simpler Frozen
ResNet-18 with Prototype Augmentation approach (47.63\% vs. 57.20\% AA). The
adapters do further reduce forgetting ($-15.54\%$ BWT), but the ResNet-18 base
features are not rich enough to anchor the task-specific adapter representations
well. This result directly motivates the move to a stronger backbone.

\subsection{Generalisation to CORe50}

We evaluated the naive baseline and our best method on CORe50 NC
\citep{lomonaco2017core50} (50 object classes, 5 tasks of 10 classes each)
to test whether the results hold beyond Split-CIFAR-100.
Table~\ref{tab:core50} shows the results alongside the CIFAR-100 numbers
for direct comparison.

\begin{table}[h]
\centering
\caption{Naive baseline vs.\ ViT + Adapters + NCM on both benchmarks.
Same class-incremental, replay-free protocol throughout.}
\label{tab:core50}
\renewcommand{\arraystretch}{1.2}
\begin{tabular}{llcc}
\toprule
\textbf{Dataset} & \textbf{Method} & \textbf{AA (\%)} & \textbf{BWT (\%)} \\
\midrule
Split-CIFAR-100 & Naive (ViT)          & 36.32 & $-69.58$ \\
Split-CIFAR-100 & ViT + Adapters + NCM & 74.81 & $-9.64$  \\
\midrule
CORe50 NC       & Naive (ViT)          & 40.11 & $-63.89$ \\
CORe50 NC       & ViT + Adapters + NCM & \textbf{86.52} & $\mathbf{-2.62}$ \\
\bottomrule
\end{tabular}
\end{table}

The adapter method improves further on CORe50 (86.52\% AA, $-2.62\%$ BWT),
likely because each task has only 10 classes instead of 20, reducing
classifier crowding. The near-zero BWT confirms that the architecture
generalises well across datasets.

\subsection{Best Result: ViT-B/16 with Per-Task Adapters and NCM}

This method brings together every design choice from the prior experiments.
A frozen ViT-B/16 backbone supplies strong, stable 768-dimensional features.
Per-task MLP adapters trained on each task and then permanently frozen add
discriminative capacity without touching anything learned before. NCM inference
over the growing concatenated feature space handles class-incremental queries
by comparing against all stored prototypes simultaneously, with no task identity
required. A Tukey power transform and L2 normalisation are applied before
distance computation.

\begin{figure}[h]
  \centering
  \includegraphics[width=0.85\linewidth]{vit_adapters_bar.png}
  \caption{Per-task accuracy after full five-task training for all four
    ViT-based methods. ViT Naive, ViT LwF, and ViT EWC all show severe
    forgetting on early tasks, with Tasks 1 and 2 dropping well below 50\%
    by the end of training. Our best method, ViT + Adapters + NCM, maintains
    consistently high accuracy across all five tasks, reaching 74.81\% AA
    and only $-9.64\%$ BWT in the class-incremental setting.}
  \label{fig:best_bar}
\end{figure}

The final result is 74.81\% AA and $-9.64\%$ BWT across five tasks
(Figure~\ref{fig:best_bar}). Task 1, which is the most exposed to subsequent
task interference, retains 75.8\% accuracy at the very end of five-task
training, down from 92.5\% immediately after its own training. Every task
achieves between 70\% and 93\% accuracy at the end of the full sequence.
Figure~\ref{fig:forgetting_vit} shows this directly: while naive fine-tuning
causes all earlier tasks to collapse as training progresses, ADAM's task lines
remain stable throughout. The single most impactful architectural decision
across all configurations is the classifier type: switching from a trainable
linear head to NCM reduces BWT from roughly $-40\%$ to roughly $-10\%$
regardless of backbone choice.

\begin{figure}[h]
  \centering
  \includegraphics[width=\linewidth]{best_result.png}
  \caption{Forgetting curves for Naive ViT (left) vs.\ ADAM (right) on
    Split-CIFAR-100. Each coloured line tracks one task's accuracy as
    subsequent tasks are trained. Under naive fine-tuning all earlier tasks
    collapse toward 0\% as soon as a new task begins. Under ADAM, every
    task's accuracy remains stable throughout training, demonstrating that
    the adapter and NCM design structurally eliminates forgetting.}
  \label{fig:forgetting_vit}
\end{figure}

% ================================================================
\section{Conclusions and Future Work}
\label{sec:conclusions}

\subsection{Main Contributions}

This project asked how much of the catastrophic forgetting problem can be solved
in a fully replay-free, class-incremental setting. The answer is: quite a lot,
provided the architecture is redesigned rather than just the training objective.

Regularization-based methods (EWC, LwF, and their combination) plateau below
22\% AA on ResNet-18 and below 42\% on ViT-B/16. No amount of hyperparameter
tuning within this family fixes the root cause: all tasks competing for the same
weights in a shared, mutable network. Regularizers slow drift but cannot stop
it.

The breakthrough comes from replacing the trainable classifier with a
non-parametric one. When there are no classification weights to overwrite,
that source of forgetting is eliminated by design. Adding per-task adapters on
top of a strong frozen backbone incrementally expands the feature space, giving
each class a richer and more discriminative description without disturbing
anything already learned. Our best model reaches 74.81\% AA and $-9.64\%$ BWT
without storing a single past training example and without requiring task
identity at test time.

\subsection{Lessons Learned}

Working through eight progressively better configurations made several things
very concrete:

\begin{enumerate}[leftmargin=*]
  \item Regularization targets a symptom (weight drift) rather than the cause
    (shared mutable capacity). It has a hard ceiling, especially in
    class-incremental settings where the classifier must simultaneously
    distinguish all accumulated classes.
  \item Non-parametric classifiers like NCM are a natural fit for replay-free
    class-incremental learning. Adding a new class prototype never overwrites
    information about an old one, and no task routing is needed at inference.
  \item Per-task adapters with permanent freezing expand representational
    capacity incrementally. The feature space grows richer over time without
    any corruption of past representations.
  \item The quality of the base backbone matters enormously. Adapters built on
    weak ResNet-18 features underperform a simple prototype-only approach. On
    ViT-B/16 features, the same adapters add substantial value.
  \item AA and BWT do not always move together. Some configurations show
    acceptable BWT but poor AA because the initial per-task accuracy was already
    low. Both metrics need to be tracked and reported together to get an honest
    picture of continual learning performance.
\end{enumerate}

\subsection{Future Work}

Several directions are worth exploring from here:

\textbf{Larger backbones.} ViT-L/16 or CLIP-pretrained models produce even
richer feature representations. Pairing those with the adapter and NCM framework
could push AA above 80\% in the class-incremental setting while keeping
forgetting near zero.

\textbf{Adapter architecture search.} Our adapters use a fixed two-layer MLP
with 128-dimensional output. Systematically exploring adapter depth, width, and
connectivity (including residual paths) could find a better tradeoff between
added capacity and inference latency.

\textbf{Harder benchmarks.} Split-CIFAR-100 has clean disjoint task boundaries.
Domain-incremental benchmarks like CORe50, or larger-scale datasets like
Split-ImageNet, would test whether these architectural choices hold up in more
realistic and challenging scenarios.

\textbf{Minimal rehearsal.} Even a very small replay buffer (around 20 examples
per class) combined with the adapter and NCM framework might close the remaining
accuracy gap without meaningfully violating the memory constraint.

\textbf{Inference scalability.} As more tasks are added, the concatenated
feature vector and prototype store both grow. At 20 or more tasks this starts
to matter for latency. Approximate nearest-neighbor search or prototype
compression could keep inference fast without hurting accuracy.

% ================================================================
\section{Code}
\label{sec:code}

All code for this project, including data preprocessing, model implementations
for all eight methods, training scripts, and evaluation utilities, is available
in our GitHub repository:

\begin{center}
  \url{https://github.com/devalladivyasreedurga/RFCL}
\end{center}

The repository is organized as follows. The \texttt{models/} directory contains
implementations of each method, from naive fine-tuning through the ViT adapter
architecture. The \texttt{train.py} script runs any of the eight methods on
Split-CIFAR-100 with configurable hyperparameters. The \texttt{evaluate.py}
script computes AA and BWT after training. The \texttt{notebooks/} directory
contains the plotting code used to generate the figures in this report.
Dependencies are listed in \texttt{requirements.txt} and the full experimental
setup is documented in the repository README.

% ================================================================
\section*{References}

\begingroup
\renewcommand{\section}[2]{}
\begin{thebibliography}{10}

\bibitem{kirkpatrick2017overcoming}
J.~Kirkpatrick, R.~Pascanu, N.~Rabinowitz, J.~Veness, G.~Desjardins,
A.~A.~Rusu, K.~Milan, J.~Quan, T.~Ramalho, A.~Grabska-Barwinska,
D.~Hassabis, C.~Clopath, D.~Kumaran, and R.~Hadsell.
\newblock Overcoming catastrophic forgetting in neural networks.
\newblock \emph{Proceedings of the National Academy of Sciences},
  114(13):3521--3526, 2017.

\bibitem{li2017learning}
Z.~Li and D.~Hoiem.
\newblock Learning without forgetting.
\newblock In \emph{European Conference on Computer Vision (ECCV)}, 2016.

\bibitem{zhang2023slca}
G.~Zhang, L.~Wang, G.~Kang, L.~Chen, and Y.~Wei.
\newblock {SLCA}: Slow learner with classifier alignment for continual learning
  on a pre-trained model.
\newblock In \emph{International Conference on Computer Vision (ICCV)}, 2023.

\bibitem{zhu2021prototype}
F.~Zhu, X.-Y.~Zhang, C.~Wang, F.~Yin, and C.-L.~Liu.
\newblock Prototype augmentation and self-supervision for incremental learning.
\newblock In \emph{Proceedings of the IEEE/CVF Conference on Computer Vision
  and Pattern Recognition (CVPR)}, pages 5871--5880, 2021.

\bibitem{zhou2024revisiting}
D.-W.~Zhou, Z.-W.~Cai, H.-J.~Ye, D.-C.~Zhan, and Z.~Liu.
\newblock Revisiting class-incremental learning with pre-trained models:
  Generalizability and adaptivity are all you need.
\newblock \emph{International Journal of Computer Vision}, 133:1012--1032,
  2024.

\bibitem{dosovitskiy2020image}
A.~Dosovitskiy, L.~Beyer, A.~Kolesnikov, D.~Weissenborn, X.~Zhai,
T.~Unterthiner, M.~Dehghani, M.~Minderer, G.~Heigold, S.~Gelly,
J.~Uszkoreit, and N.~Houlsby.
\newblock An image is worth 16x16 words: Transformers for image recognition
  at scale.
\newblock In \emph{International Conference on Learning Representations
  (ICLR)}, 2021.

\bibitem{lomonaco2017core50}
V.~Lomonaco and D.~Maltoni.
\newblock {CORe50}: a new dataset and benchmark for continuous object
  recognition.
\newblock In \emph{Conference on Robot Learning (CoRL)}, 2017.

\bibitem{zenke2017continual}
F.~Zenke, B.~Poole, and S.~Ganguli.
\newblock Continual learning through synaptic intelligence.
\newblock In \emph{International Conference on Machine Learning (ICML)}, 2017.

\end{thebibliography}
\endgroup

\end{document}